% ============================================================
%  MODULE III — Market Participants & Structure
% ============================================================

\chapter{Market Participants \& Structure}

\begin{notebox}
Understanding \emph{who} is in the market and \emph{why} they trade changes everything.
Price action is the aggregate footprint of all participants acting on their different
objectives, time horizons, and information sets.
\end{notebox}

% ────────────────────────────────────────────────────────────
\section{The Buy Side vs. The Sell Side}

\begin{keyterm}{Sell Side}
Firms that \emph{create and sell} financial products, execute client trades, and provide
research. They earn through commissions, spreads, and underwriting fees.\\
Examples: Goldman Sachs, Morgan Stanley, JPMorgan (as investment banks), market makers.
\end{keyterm}

\begin{keyterm}{Buy Side}
Firms that \emph{buy} securities with the goal of generating returns for themselves or clients.\\
Examples: mutual funds, hedge funds, pension funds, insurance companies, endowments.
\end{keyterm}

% ────────────────────────────────────────────────────────────
\section{Types of Market Participants}

\subsection{Retail Investors}
Individual non-professional investors. Trade through brokers (Fidelity, Schwab, Interactive
Brokers, Robinhood). Generally small position sizes. Collectively significant but individually
price-takers.

\subsection{Institutional Investors}
Manage large pools of capital:
\begin{itemize}
  \item \textbf{Mutual Funds} — pool retail money, must disclose holdings quarterly, cannot short.
  \item \textbf{Pension Funds} — very long-duration mandates; large in bonds and equities.
  \item \textbf{Insurance Companies} — must match asset duration to liability duration; bond-heavy.
  \item \textbf{Endowments} — universities, foundations; long-horizon, often hold alternatives.
  \item \textbf{Sovereign Wealth Funds} — government-owned, very large (e.g., Norway GPFG, GIC).
\end{itemize}

\subsection{Hedge Funds}
Lightly regulated investment vehicles for sophisticated investors. Can use:
leverage, short selling, derivatives, concentrated positions. Performance fee structure:
typically \textbf{2 and 20} (2\% management fee + 20\% of profits).

Strategies include: long/short equity, global macro, statistical arbitrage, event-driven,
distressed credit, systematic/quant.

\subsection{Market Makers}
Provide continuous two-sided quotes (bid and ask). Profit from the \textbf{bid-ask spread}.
Obligated to buy and sell; absorb order imbalances. Essential for market liquidity.
Examples: Citadel Securities, Virtu Financial, Jane Street.

\subsection{High-Frequency Traders (HFT)}
Use ultra-fast execution (microseconds) and algorithmic strategies to profit from tiny,
fleeting price inefficiencies. Controversial — seen by some as liquidity providers, by
others as front-runners.

\subsection{Central Banks}
The Federal Reserve (US), ECB (Europe), BOJ (Japan) can buy and sell assets directly.
Quantitative easing (QE) involves large-scale asset purchases, injecting liquidity.
Central bank actions have outsized market impact.

% ────────────────────────────────────────────────────────────
\section{How a Trade Is Executed}

The lifecycle of a typical equity trade:

\begin{enumerate}[label=\textbf{Step \arabic*.}]
  \item \textbf{Order Entry} — Investor places order via broker (market, limit, stop, etc.)
  \item \textbf{Routing} — Broker routes to exchange, ATS (Alternative Trading System), or
        internalises against market maker flow.
  \item \textbf{Matching} — Order matching engine pairs buy and sell orders at a price.
  \item \textbf{Execution} — Trade confirmed. Price and quantity locked in.
  \item \textbf{Clearing} — Clearinghouse (e.g., DTCC in the US) becomes the counterparty to
        both sides, managing counterparty risk.
  \item \textbf{Settlement} — Actual transfer of securities and cash.
        US equities: \textbf{T+1} (trade date plus one business day).
\end{enumerate}

% ────────────────────────────────────────────────────────────
\section{Order Types}

\begin{center}
\begin{tabular}{>{\bfseries}p{3.5cm} p{9cm}}
\toprule
Order Type & Description \\
\midrule
Market Order & Execute immediately at best available price. Fast but no price guarantee. \\
Limit Order & Execute only at a specified price or better. Price control but may not fill. \\
Stop Order & Becomes a market order once a trigger price is hit. Used to cut losses. \\
Stop-Limit & Becomes a \emph{limit} order once trigger is hit. Risk of non-fill in fast markets. \\
Trailing Stop & Stop that moves with the price, locking in gains. \\
Iceberg Order & Only a small portion of total size is visible; remainder is hidden. \\
Fill or Kill & Execute the entire order immediately or cancel it entirely. \\
\bottomrule
\end{tabular}
\end{center}

\begin{warningbox}
In fast-moving, illiquid, or gap-prone markets, \textbf{stop orders} can fill significantly
worse than the trigger price (\emph{slippage}). Never assume a stop order equals a
guaranteed exit price.
\end{warningbox}

% ────────────────────────────────────────────────────────────
\section{Brokers and Execution Venues}

\subsection{Types of Brokers}
\begin{itemize}
  \item \textbf{Full-service broker} — advice, research, wealth management. Higher fees.
  \item \textbf{Discount broker} — execution only, low/zero commission. (Schwab, Fidelity, IBKR)
  \item \textbf{Direct Access Broker} — routes directly to exchanges; used by active traders and HFT.
        (Interactive Brokers Pro, Lightspeed, DAS Trader)
\end{itemize}

\subsection{Payment for Order Flow (PFOF)}
Retail brokers like Robinhood route orders to market makers (e.g., Citadel Securities) in
exchange for payment. Controversial: critics argue retail gets slightly worse execution;
defenders say it funds zero-commission trading. \textbf{Banned in the UK and EU.}

\subsection{Dark Pools}
Private exchanges where large institutional orders are matched anonymously, away from public
markets, to minimise \textbf{market impact}. A large sell order on a public exchange signals
intent and can move the price adversely before execution completes.
